% =====================================================================
\chapter{Búsqueda evolutiva de programas guiada por LLM}
\label{tema:1}
% =====================================================================

En muchos problemas lo que se busca no es un número ni un vector, sino un \emph{procedimiento}: una regla de prioridad para un planificador, una heurística de empaquetado, un núcleo de cálculo más rápido. Este tema estudia cómo un modelo de lenguaje, encerrado en un bucle evolutivo con un evaluador automático, puede buscar en el espacio de programas. Sigue el artículo de FunSearch de Romera-Paredes \emph{et al.}~\cite{funsearch} y su implementación de referencia~\cite{funsearch-code}, el informe de AlphaEvolve de Novikov \emph{et al.}~\cite{alphaevolve} y la nota técnica de Google DeepMind~\cite{ae-blog}; para la representación de programas como expresiones se apoya en el capítulo «Expression Optimization» de Kochenderfer y Wheeler~\cite{kw}, y para el empaquetado, en la §3.3 de Williamson y Shmoys~\cite{ws}. La lectura crítica de Herrmann y Pallez~\cite{hp} sirve de contrapeso.

\section{Optimizar en el espacio de programas}

\begin{definicion}[Búsqueda de programas]\label{def:busqueda-programas}
Sea $\Prog$ el conjunto de programas que respetan una firma dada (por ejemplo, funciones que reciben un objeto y la lista de contenedores abiertos y devuelven una puntuación por contenedor). Sea $\mathcal{D}$ una distribución de instancias del problema y $u(\pi,x)\in\R\cup\{-\infty\}$ la utilidad de ejecutar $\pi$ sobre la instancia $x$, con $u=-\infty$ si el programa falla, excede el tiempo o produce una salida inválida. La \emph{búsqueda de programas} es
\[
\max_{\pi\in\Prog}\ h(\pi),\qquad h(\pi)=\E_{x\sim\mathcal{D}}\bigl[u(\pi,x)\bigr].
\]
Como $\mathcal{D}$ no se conoce, se usa un \emph{evaluador} $\hat h(\pi)=\frac1m\sum_{i=1}^{m}u(\pi,x_i)$ sobre instancias de entrenamiento $x_1,\dots,x_m$, y la calidad final se mide con instancias de prueba independientes.
\end{definicion}

El problema es de caja negra en el sentido de la definición~5.1 del módulo~24: no hay gradiente, el espacio es discreto y enorme, y cada evaluación cuesta una ejecución. La programación genética clásica busca en ese espacio con mutaciones y cruces aleatorios de árboles sintácticos~\cite{kw}; su dificultad es que casi todas las mutaciones aleatorias producen programas inútiles. La idea de FunSearch es sustituir la mutación aleatoria por un modelo de lenguaje entrenado con código, que propone modificaciones \emph{plausibles}.

\begin{observacion}[Por qué buscar programas y no soluciones]\label{obs:por-que-programas}
(i) Un programa \emph{generaliza}: una heurística de empaquetado sirve para instancias nuevas, mientras que un empaquetado concreto solo sirve para la suya. (ii) Un programa corto puede \emph{describir} un objeto enorme: en el problema del \emph{cap set}, FunSearch buscó programas que construyen el conjunto, no el conjunto~\cite{funsearch}. (iii) Un programa se puede \emph{leer}: el experto puede entender, simplificar y auditar la regla encontrada. (iv) Un programa se puede \emph{verificar}: el evaluador lo ejecuta y comprueba la validez de la salida.
\end{observacion}

\begin{definicion}[Esqueleto y función que evoluciona]\label{def:esqueleto}
Un \emph{esqueleto} es un programa fijo que resuelve el problema llamando a una función con firma conocida, la única que se modifica durante la búsqueda. En el empaquetado en línea, el esqueleto recorre los objetos en orden de llegada, calcula los contenedores en los que cabe cada uno y coloca el objeto en el de mayor \texttt{prioridad(objeto, contenedores)}; solo \texttt{prioridad} evoluciona.
\end{definicion}

El esqueleto encierra el conocimiento del problema y garantiza la factibilidad (un objeto nunca se coloca donde no cabe), de modo que cualquier función que devuelva números produce una solución válida. Es el mismo principio de diseño que el filtro de seguridad del módulo~26 (tema~6, §6.5): la parte aprendida propone dentro de un marco que impide lo inaceptable.

\section{FunSearch}

\addtocounter{figure}{1}%

El bucle de la \cref{fig:funsearch} se ejecuta de forma asíncrona y distribuida. Según el artículo y la implementación publicada~\cite{funsearch,funsearch-code}:
\begin{enumerate}[label=\arabic*.]
  \item \emph{Constructor del prompt.} Se elige una isla al azar y, dentro de ella, $k=2$ programas; se ordenan de peor a mejor puntuación, se renombran como \texttt{prioridad\_v0} y \texttt{prioridad\_v1}, y el prompt termina con la cabecera de \texttt{prioridad\_v2}, «versión mejorada de la v1». A esta técnica los autores la llaman \emph{best-shot prompting}.
  \item \emph{Muestreo.} Un LLM preentrenado con código (en el trabajo original, un modelo de la familia PaLM~2; en la actualización de 2024, Gemini~1.5 Flash~\cite{fs-blog}) completa el cuerpo de la nueva función. La configuración por defecto de la implementación usa 15 muestreadores que piden 4 continuaciones por prompt.
  \item \emph{Evaluación.} El programa se ejecuta en un entorno aislado sobre las instancias de entrenamiento; si falla, tarda demasiado o produce algo inválido, se descarta. La configuración por defecto usa 140 evaluadores en paralelo.
  \item \emph{Base de programas.} Los programas válidos se guardan en $10$ islas que evolucionan por separado. Dentro de cada isla se agrupan en \emph{clusters} por su \emph{firma} (el vector de puntuaciones en cada instancia de entrenamiento). Periódicamente se reinicia la mitad de las islas con peor mejor puntuación, sembrando cada una con el mejor programa de una isla superviviente elegida al azar.
\end{enumerate}

\begin{definicion}[Muestreo de clusters con temperatura]\label{def:softmax-clusters}
Sean $s_1,\dots,s_c$ las puntuaciones de los clusters de una isla y $T>0$ una temperatura. Se elige el cluster $i$ con probabilidad
\[
p_i(T)=\frac{\exp(s_i/T)}{\sum_{l=1}^{c}\exp(s_l/T)}.
\]
En la implementación de referencia, $T$ decrece linealmente con el número $n$ de programas registrados en la isla, $T=T_0\bigl(1-(n\bmod P)/P\bigr)$, con $T_0=0{,}1$ y periodo $P=30\,000$; dentro del cluster elegido se prefieren los programas más cortos, con una softmax sobre la longitud normalizada cambiada de signo~\cite{funsearch-code}.
\end{definicion}

\begin{proposicion}[Límites de la temperatura]\label{prop:temperatura}
Sea $M=\max_i s_i$ y $A=\{i: s_i=M\}$. Cuando $T\to 0^{+}$, $p_i(T)\to 1/\abs{A}$ si $i\in A$ y $p_i(T)\to 0$ en otro caso. Cuando $T\to\infty$, $p_i(T)\to 1/c$ para todo $i$.
\end{proposicion}

La temperatura regula, pues, el equilibrio entre \emph{explotar} los mejores programas y \emph{explorar} otros que pueden ser punto de partida de mejoras mayores; es la misma distribución de Boltzmann del recocido simulado (módulo~24, tema~5, §5.2). Las islas y los reinicios cumplen una función parecida a otra escala: mantienen poblaciones que no se contaminan entre sí y, de vez en cuando, abandonan las que se han estancado.

\begin{proposicion}[Elitismo y espera de una mejora]\label{prop:espera}
Sea $b_t$ la mejor puntuación de la base de programas tras $t$ evaluaciones, en un algoritmo que nunca elimina el mejor programa. Entonces:
\begin{enumerate}
  \item $b_{t+1}\ge b_t$ para todo $t$.
  \item Si, mientras no haya mejora, cada propuesta supera a $b_t$ con probabilidad al menos $q>0$, de forma independiente, el número $N$ de evaluaciones hasta la primera mejora cumple $\E[N]\le 1/q$ y $\Prob(N>B)\le(1-q)^{B}\le e^{-qB}$.
  \item Si $I$ islas independientes tienen cada una probabilidad $p$ de alcanzar una puntuación objetivo con su presupuesto, la probabilidad de que la alcance alguna es $1-(1-p)^{I}$.
\end{enumerate}
\end{proposicion}

La proposición explica el papel del LLM con una sola cifra: el presupuesto necesario es del orden de $1/q$. Una mutación aleatoria de un árbol sintáctico tiene una $q$ minúscula; un modelo que ha visto millones de programas concentra la probabilidad en ediciones que tienen sentido (cambiar un umbral, añadir un caso, simplificar una expresión), y eso multiplica $q$.

\section{AlphaEvolve}

AlphaEvolve~\cite{alphaevolve} generaliza FunSearch de función a sistema. El usuario marca con comentarios (\texttt{EVOLVE-BLOCK-START} y \texttt{EVOLVE-BLOCK-END}) los bloques de un fichero que pueden cambiar; el resto actúa como esqueleto. Los LLM devuelven \emph{diferencias} en formato de búsqueda y sustitución en lugar de reescribir el fichero, lo que permite evolucionar programas de cientos de líneas en cualquier lenguaje. Un conjunto de modelos combina uno rápido (Gemini~2.0 Flash, muchas propuestas) y uno más capaz (Gemini~2.0 Pro, menos propuestas y de más calidad). La evaluación es una \emph{cascada}: el programa pasa por pruebas de dificultad creciente y solo avanza si supera las anteriores, de modo que los candidatos malos se descartan con poco coste. La base de programas combina MAP-Elites (una población que guarda el mejor programa en cada celda de un espacio de características, para mantener diversidad) con el modelo de islas. El \cref{tab:fs-ae} resume las diferencias que los propios autores señalan.

\begin{table}[h]
\centering\small
\caption{FunSearch frente a AlphaEvolve, según la tabla~1 de~\cite{alphaevolve}.}\label{tab:fs-ae}
\begin{tabularx}{\textwidth}{@{}>{\raggedright\arraybackslash}p{3.0cm}>{\raggedright\arraybackslash}X>{\raggedright\arraybackslash}X@{}}
\toprule
 & FunSearch & AlphaEvolve \\
\midrule
qué evoluciona & una función & un fichero entero, con varios bloques \\
tamaño & 10--20 líneas & hasta cientos de líneas \\
lenguaje & Python & cualquiera \\
evaluación & $\le 20$~min en una CPU & horas, en paralelo, con aceleradores \\
muestras del LLM & millones & bastan miles \\
\bottomrule
\end{tabularx}
\end{table}

\addtocounter{definicion}{1}%

\section{Por qué funciona y cuándo no}

Tres ingredientes explican el éxito. \emph{Asimetría entre generar y verificar}: comprobar que un empaquetado es válido y contar contenedores es inmediato, aunque encontrar una buena heurística sea difícil; es la misma asimetría que explota el aprendizaje por refuerzo con recompensas verificables (módulo~21, tema~6, §6.2) y el esquema de generar y verificar del módulo~27 (tema~6, §6.3). \emph{Una distribución de propuestas informada}: el LLM eleva la probabilidad $q$ de la \cref{prop:espera}. \emph{Memoria y diversidad}: la base de programas acumula lo aprendido y las islas evitan la convergencia prematura.

Los límites son igual de claros. (i) \emph{El evaluador es la especificación}: lo que no mida no se optimiza, y lo que mida mal se explota, como en la sobreoptimización de la recompensa (módulo~21, tema~6, §6.4). (ii) \emph{Generalización}: la heurística se ajusta a la distribución de entrenamiento; Herrmann y Pallez~\cite{hp} concluyen que las heurísticas de empaquetado de FunSearch no son robustas a cambios en la distribución de los datos y que en gran parte redescubren patrones conocidos. (iii) \emph{Coste}: millones de muestras en FunSearch y horas de evaluación en AlphaEvolve. (iv) \emph{Seguridad}: el código generado es código no confiable y debe ejecutarse aislado, sin red ni acceso a datos sensibles.

\section{Empaquetado en línea}

\begin{definicion}[Empaquetado unidimensional en línea]\label{def:bp}
Se reciben, uno a uno, objetos de tamaños $s_1,s_2,\dots,s_N\in(0,C]$ que deben colocarse en contenedores de capacidad $C$ sin superarla. Cada objeto se coloca al llegar, sin conocer los siguientes, y no se puede mover después. El objetivo es minimizar el número de contenedores usados. Sea $\mathrm{OPT}$ el mínimo que se lograría conociendo todos los objetos de antemano.
\end{definicion}

\begin{definicion}[Reglas de prioridad clásicas]\label{def:ff-bf}
Sea $r$ la capacidad restante de un contenedor abierto en el que cabe el objeto $s$ y $g=r-s$ el hueco que dejaría. \emph{First fit} (FF) coloca el objeto en el primer contenedor abierto en que cabe; \emph{best fit} (BF), en el que deja menor hueco $g$; en ambos casos se abre uno nuevo solo si no cabe en ninguno. En términos del esqueleto, FF usa la prioridad $-(\text{índice del contenedor})$ y BF la prioridad $-g$. Un algoritmo es de tipo \emph{any fit} si nunca abre un contenedor nuevo cuando el objeto cabe en uno abierto.
\end{definicion}

\begin{lema}[Cota $L_1$]\label{lem:l1}
$\mathrm{OPT}\ge L_1=\bigl\lceil\sum_{j}s_j/C\bigr\rceil$.
\end{lema}

\begin{teorema}[Garantía de los algoritmos \emph{any fit}]\label{teo:anyfit}
Todo algoritmo \emph{any fit}, en particular FF y BF, usa a lo sumo $\lceil 2\sum_j s_j/C\rceil\le 2\,\mathrm{OPT}$ contenedores.
\end{teorema}

La constante~2 no es la mejor posible para FF, pero el resultado ilustra el tipo de garantía que una heurística generada \emph{no} trae consigo: lo que se sabe de ella es lo que mide el evaluador. Para el problema sin restricciones de llegada, Williamson y Shmoys demuestran que, salvo que $\mathrm{P}=\mathrm{NP}$, ningún algoritmo de tiempo polinómico garantiza un factor menor que $3/2$~\cite[§3.3]{ws}.

\addtocounter{definicion}{1}%

\section{Caso de estudio: una heurística de empaquetado generada}

\addtocounter{definicion}{1}%

\begin{table}[h]
\centering\small
\caption{Empaquetado en línea de 5 instancias de prueba de 5000 objetos ($\bar L_1=1985{,}6$). Datos del cuaderno.}\label{tab:bp}
\begin{tabular}{@{}lcc@{}}
\toprule
Regla & Bandejas (media) & Exceso sobre $L_1$ \\
\midrule
\emph{first fit} & 2073{,}4 & 4{,}42\,\% \\
\emph{best fit} & 2064{,}8 & 3{,}99\,\% \\
regla evolucionada $\pi^{*}$ & 2011{,}2 & 1{,}29\,\% \\
heurística publicada por FunSearch & 2000{,}0 & 0{,}73\,\% \\
\bottomrule
\end{tabular}
\end{table}

Como referencia externa, en el conjunto Weibull de 5000 objetos que acompaña al artículo, BF usa de media 2067{,}0 bandejas y la heurística de FunSearch 2001{,}4, con $\bar L_1=1987{,}8$: excesos del $3{,}98\,\%$ y del $0{,}68\,\%$~\cite{funsearch-code}. La \emph{prueba de robustez} del cuaderno repite la comparación con otra distribución, Weibull de escala 60 (objetos más grandes): BF deja un exceso del $5{,}19\,\%$, la regla $\pi^{*}$ un $5{,}21\,\%$ y la heurística publicada un $6{,}40\,\%$, peor que FF ($6{,}05\,\%$). Con seis semillas distintas del generador, el exceso en prueba de la regla encontrada varía entre el $1{,}29\,\%$ y el $2{,}33\,\%$, y cuatro de las seis reglas tienen la misma estructura: un término de \emph{best fit} saturado más una penalización de los huecos intermedios. Las dos heurísticas generadas se han especializado en la distribución con la que se evaluaron, tal como advierten Herrmann y Pallez~\cite{hp}.

\begin{table}[h]
\centering\small
\caption{Tecnologías del tema~1.}\label{tab:tec1}
\begin{tabularx}{\textwidth}{@{}>{\raggedright\arraybackslash}p{2.5cm}>{\raggedright\arraybackslash}X>{\raggedright\arraybackslash}p{2.1cm}>{\raggedright\arraybackslash}X@{}}
\toprule
Herramienta & Para qué sirve & Licencia o tipo & Cuándo elegirla \\
\midrule
FunSearch (repositorio)~\cite{funsearch-code} & algoritmo evolutivo, manipulación de código y una versión de un solo hilo, sin LLM ni entorno aislado & abierta (Apache~2.0) & entender el método y adaptarlo a un LLM y una infraestructura propios \\
OpenEvolve~\cite{openevolve} & implementación abierta inspirada en AlphaEvolve: islas, MAP-Elites, evaluación en cascada, conjuntos de LLM & abierta (Apache~2.0) & evolucionar ficheros completos con cualquier API compatible \\
AlphaEvolve~\cite{alphaevolve} & agente de Google DeepMind & servicio, no publicado como código & referencia de diseño; resultados publicados \\
gplearn, DEAP & programación genética clásica sin LLM & abiertas (BSD-3; LGPL) & línea base barata para expresiones pequeñas \\
Contenedores aislados & ejecutar código no confiable sin red y con límites de tiempo y memoria & abierta o del proveedor & siempre que se ejecute código generado \\
\bottomrule
\end{tabularx}
\end{table}

\addtocounter{definicion}{1}%

